\subsection{FUNDAMENTAL SYSTEMS OF INTEGRALS OF A LINEAR SYSTEM OF SCALAR DIFFERENTIAL EQUATIONS}

In this subsection and the next we shall consider the case where E is a vector space of finite dimension n over the field C of complex numbers (so of dimension 2n over R), and where, for every \(t \in J\) , \(A(t)\) is an endomorphism of E with respect to the vector space structure over C. One can then identify \(A(t)\) with its matrix \((a_{ij}(t))\) with respect to a basis of E (over the field C), the \(a_{ij}\) this time being \(n^{2}\) complex functions defined and regulated on J; letting \(x_{j}\) ( \(1 \leqslant j \leqslant n\) ) denote the (complex) components of a vector \(x \in E\) with respect to the chosen basis, the linear equation

\[
{\frac {d \mathbf {x}}{d t}} = A (t). \mathbf {x} + \mathbf {b} (t)\tag{2}
\]

is again equivalent to the system

\[
\frac {d x _ {i}}{d t} = \sum_ {j = 1} ^ {n} a _ {i j} (t) x _ {j} + b _ {i} (t) \quad (1 \leqslant i \leqslant n).\tag{3}
\]

Theorems 1 (IV, p.~179) and 2 (IV, p.~179) and prop. 2 (IV, p.~181) then show that for every \(\mathbf{x}_0 = (x_{k0})_{1\leqslant k\leqslant n}\) in E there exists one and only one integral \(\mathbf{u} = (u_k)_{1\leqslant k\leqslant n}\) of the equation

\[
{\frac {d \mathbf {x}}{d t}} = A (t). \mathbf {x}\tag{4}
\]

defined on E and equal to \(x_{0}\) at the point \(t_{0}\) ; this integral can be written

\[
\mathbf {u} (t, t _ {0}, \mathbf {x} _ {0}) = C (t, t _ {0}). \mathbf {x} _ {0},
\]

\(C(t, t_{0})\) being an invertible square matrix \((c_{ij}(t, t_{0}))\) of order n whose entries are continuous complex functions on J × J and such that \(t \mapsto c_{ij}(t, t_{0})\) is a primitive of a regulated function on J.

In the particular case where \(n = 1\) the system (3) reduces to a single scalar equation

\[
{\frac {d x}{d t}} = a (t) x + b (t)\tag{11}
\]

\((a(t) \text{ and } b(t) \text{ being complex regulated functions on J}); \text{ one verifies immediately that the (one element) matrix } C(t, t_0) \text{ is equal to } \exp\left(\int_{t_0}^{t} a(s) ds\right); \text{ the integral of (11) equal to } x_0 \text{ at the point } t_0 \text{ is thus given explicitly by the formula}\)

\[
u (t) = x _ {0} \exp \left(\int_ {t _ {0}} ^ {t} a (s) d s\right) + \int_ {t _ {0}} ^ {t} b (s) \exp \left(\int_ {t _ {0}} ^ {t} a (\tau) d \tau\right) d s.\tag{12}
\]

In the space \(\mathcal{C}(\mathrm{J};\mathrm{E})\) of continuous maps from J into E, endowed with the topology of compact convergence, the set I of integrals of equation (4) is a vector subspace (over C) isomorphic to E, therefore to \(C^{n}\) (IV, p.~180, cor. 1, and IV, p.~181, prop. 2).

A basis \((\mathbf{u}_{j})_{1\leqslant j\leqslant n}\) of this space (over the field C) is called a fundamental system of integrals of (4).

PROPOSITION 4. For the n integrals \(\mathbf{u}_{j}\ (1 \leqslant j \leqslant n)\) of equation (4) to form a fundamental system it is necessary and sufficient that their values \(\mathbf{u}_{j}(t_{0})\) at a point \(t_{0} \in J\) be linearly independent vectors in E.

Indeed, the map which to every \(\mathbf{x}_0 \in \mathrm{E}\) associates the integral \(t \mapsto C(t, t_0). \mathbf{x}_0\) is an isomorphism of \(\mathrm{E}\) onto \(\mathcal{I}\) (IV, p.~180, cor. 1 and IV, p.~181, prop. 2).

If \(\left(\mathbf{e}_j\right)_{1\leqslant j\leqslant n}\) is any basis of E over C, the \(n\) integrals

\[
\mathbf {u} _ {j} (t) = C \left(t, t _ {0}\right). \mathbf {e} _ {j} \quad (1 \leqslant j \leqslant n)
\]

thus form a fundamental system; if one identifies \(C(t, t_0)\) with its matrix with respect to the basis \((\mathbf{e}_j)\) the integrals \(\mathbf{u}_j\) are precisely the columns of the matrix \(C(t, t_0)\) . The integral of (4) that takes the value \(\mathbf{x}_0 = \sum_{j=1}^{n} \lambda_j \mathbf{e}_j\) at the point \(t_0\) is then \(C(t, t_0). \mathbf{x}_0 = \sum_{k=1}^{n} \lambda_k \mathbf{u}_k(t)\) .

Given any \(n\) integrals \(\mathbf{u}_j\) ( \(1 \leqslant j \leqslant n\) ) of (4) one terms the determinant of these \(n\) integrals at a point \(t \in \mathbf{J}\) with respect to a basis \((\mathbf{e}_j)_{1 \leqslant j \leqslant n}\) of \(E\) , the determinant

\[
\Delta (t) = \left(\mathbf {u} _ {1} (t), \mathbf {u} _ {2} (t), \dots , \mathbf {u} _ {n} (t)\right)\tag{13}
\]

of the \(n\) vectors \(\mathbf{u}_j(t)\) with respect to the basis \((\mathbf{e}_j)\) (Alg., III, p.~522). One has (Alg., III, p.~523, prop. 2)

\[
\Delta (t) = \Delta (t _ {0}) \det \left(C (t, t _ {0})\right).\tag{14}
\]

By prop. 4 of IV, p.~184, for \(\left(\mathbf{u}_{j}\right)_{1\leqslant j\leqslant n}\) to be a fundamental system of integrals of (4) it is necessary and sufficient that the determinant \(\Delta(t)\) of the \(u_{j}\) be \(\neq0\) at some one point \(t_{0}\) of J; the formula (14) then shows that \(\Delta(t)\neq0\) at every point of J, in other words, that the vectors \(\mathbf{u}_{j}(t)\) ( \(1\leqslant j\leqslant n\) ) are always linearly independent.

PROPOSITION 5. The determinant of the matrix \(C(t, t_{0})\) is given by the formula

\[
\det \left(C (t, t _ {0})\right) = \exp \left(\int_ {t _ {0}} ^ {t} \operatorname{Tr} (A (s)) d s\right).\tag{15}
\]

Indeed, if one puts \(\delta(t) = \det\left(C(t, t_0)\right)\) one has, by the formula for the derivative of a determinant (I, p.~8, formula (3))

\[
\frac {d \delta}{d t} = \operatorname{Tr} \left(\frac {d C (t , t _ {0})}{d t} \left(C (t, t _ {0})\right) ^ {- 1}\right) \delta (t)
\]

that is, by the differential equation (5) of IV, p.~179 satisfied by \(C(t, t_0)\) ,

\[
\frac {d \delta}{d t} = \mathrm{Tr} \big (A (t) \big) \delta (t).
\]

Since \(\delta(t_0) = 1\) the formula (15) follows from the expression (12) (IV, p.~183) for the integral of a scalar linear equation.

Specifying n linearly independent integrals of (4) determines all the integrals of this equation, as we have just seen. We shall now show that for \(1 \leqslant p \leqslant n\) the knowledge of p linearly independent integrals \(\mathbf{u}_{j}\) ( \(1 \leqslant j \leqslant p\) ) of equation (4) reduces the integration of this equation to that of a homogeneous linear system of n - p scalar equations. Suppose that on an interval \(K \subset J\) there are n - p maps

\[
\mathbf {u} _ {p + k} \quad (1 \leqslant k \leqslant n - p)
\]

of K into E, which are primitives of regulated functions on K, and such that, for every \(t \in K\) , the n vectors \(\mathbf{u}_{j}(t) (1 \leqslant j \leqslant n)\) form a basis of E.

For every point \(t_1 \in \mathbf{J}\) there always exists an interval \(\mathbf{K}\) , a neighbourhood of \(t_1\) in \(\mathbf{J}\) , on which there are defined \(n - p\) functions \(\mathbf{u}_{p+k}\) ( \(1 \leqslant k \leqslant n - p\) ) having the preceding properties. For, let \((\mathbf{e}_i)_{1 \leqslant i \leqslant n}\) be a basis of \(\mathbf{E}\) ; there exist \(n - p\) vectors of this basis which form with the \(\mathbf{u}_j(t_1)\) ( \(1 \leqslant j \leqslant p\) ) a basis of \(\mathbf{E}\) (Alg., II, p.~292, th. 2); suppose for example that they are \(\mathbf{e}_{p+1}, \ldots, \mathbf{e}_n\) ; since the determinant \(\det(\mathbf{u}_1(t), \ldots, \mathbf{u}_p(t), \mathbf{e}_{p+1}, \ldots, \mathbf{e}_n)\) (with respect to the basis \((\mathbf{e}_i)\) ) is a continuous function of \(t\) and does not vanish for \(t = t_1\) there exists a neighbourhood \(\mathbf{K}\) of \(t_1\) on which it does not vanish; one can then take \(\mathbf{u}_{p+k}(t) = \mathbf{e}_{p+k}\) ( \(1 \leqslant k \leqslant n - p\) ) for \(t \in \mathbf{K}\) .

There exists an invertible matrix \(B(t)\) of order n, whose elements are primitives of regulated functions on K, such that \(B(t).e_{j} = \mathbf{u}_{j}(t)\) for \(1 \leqslant j \leqslant n\) . Let us put \(\mathbf{x} = B(t).y\) ; then y satisfies the equation \(\frac{dB}{dt}.y + B(t).\frac{dy}{dt} = A(t)B(t).y\) , which can also be written

\[
\frac {d \mathbf {y}}{d t} = (B (t)) ^ {- 1} \left(A (t) B (t) - \frac {d B}{d t}\right). \mathbf {y} = H (t). \mathbf {y}
\]

where \(H(t) = \left(h_{jk}(t)\right)\) is a matrix with regulated entries on K. By the definition of \(B(t)\) this linear equation admits the p constant vectors \(\mathbf{e}_{j}\) ( \(1 \leqslant j \leqslant p\) ) as integrals; one concludes immediately that necessarily \(h_{jk}(t) = 0\) for \(1 \leqslant k \leqslant p\) ; thus the components \(y_{k}\) of y (with respect to the basis \((\mathbf{e}_{i})\) ) with index \(k \geqslant p + 1\) satisfy a homogeneous linear system of n - p equations; once the solutions of this system are determined, the \(dy_{j}/dt\) for indices \(j \leqslant p\) are linear functions of the \(y_{k}\) with \(k \geqslant p + 1\) , so are known, and the primitives of these functions will give the \(y_{j}\) for indices \(j \leqslant p\) .

In particular, when one knows \(n - 1\) linearly independent integrals of equation (4) of IV, p.~183, integrating this equation reduces to that of a single homogeneous scalar equation, and then the evaluation of \(n\) primitives.

Remarks. 1) All the above applies also to the case where E is of dimension n over the field R and \(A(t)\) is an endomorphism of E for every \(t \in J\) : one has only to replace C by R throughout.

\begin{enumerate}
\def\labelenumi{\arabic{enumi})}
\setcounter{enumi}{1}
\tightlist
\item
  Let \(A(t) = (a_{ij}(t))\) be a square matrix of order n whose entries are real (resp. complex) regulated functions of t on J, and let \(C(t, t_0) = (c_{ij}(t, t_0))\) be the resolvent matrix of the corresponding linear system (3) (IV, p.~22). Let F be an arbitrary complete normed space over R (resp. C) and let us consider the system of linear differential equations
\end{enumerate}

\[
\frac {d \mathbf {y} _ {i}}{d t} = \sum_ {j = 1} ^ {n} a _ {i j} (t) \mathbf {y} _ {j}
\]

where the unknown functions \(\mathbf{y}_j\) take their values in F. It is immediate that the solution \((\mathbf{u}_j)_{1\leqslant j\leqslant n}\) of this system such that \(\mathbf{u}_j(t_0) = \mathbf{d}_j\) for \(1\leqslant j\leqslant n\) ( \(\mathbf{d}_j\) being arbitrary in F) is given by the formulae

\[
\mathbf {u} _ {i} (t) = \sum_ {j = 1} ^ {n} c _ {i j} (t, t _ {0}) \mathbf {d} _ {j} \quad (1 \leqslant i \leqslant n).
\]

We consider in particular the case where \(A(t)\) is an endomorphism of a vector space E of finite dimension n over C, such that there exists a basis of E with respect to which the matrix of \(A(t)\) has real elements for every \(t \in J\) . Then the above shows (by th. 1 of IV, p.~179) that the resolvent matrix \(C(t, t_{0})\) with respect to the same basis also has real elements: it suffices to consider the vector space \(E_{0}\) over R generated by the basis of E under consideration, and to remark that the restriction of \(A(t)\) to \(E_{0}\) is an endomorphism of this vector space.

